\documentclass[12pt,a4paper]{article}
% Package dan pengaturan umum. Konfigurasi/cover ada di main.tex.
\usepackage[utf8]{inputenc}
\usepackage[T1]{fontenc}
\usepackage[indonesian]{babel}
\usepackage[margin=2.5cm]{geometry}
\usepackage{graphicx}
\usepackage{booktabs}
\usepackage{longtable}
\usepackage{array}
\usepackage{amsmath,amssymb}
\usepackage{listings}
\usepackage{xcolor}
\usepackage{float}
\usepackage{caption}
\usepackage{subcaption}
\usepackage{tikz}
\usepackage{setspace}
\usepackage{hyperref}
\usepackage{fancyhdr}

\graphicspath{{figures/}}
\setstretch{1.15}
\hypersetup{colorlinks=true, linkcolor=black, urlcolor=blue, citecolor=black}

\lstset{
  basicstyle=\ttfamily\small,
  breaklines=true,
  frame=single,
  numbers=left,
  numberstyle=\tiny,
  keywordstyle=\color{blue},
  commentstyle=\color{gray},
}

\pagestyle{fancy}
\fancyhf{}
\fancyhead[L]{IF4020 Kriptografi}
\fancyhead[R]{Tugas 2: \algname}
\fancyfoot[C]{\thepage}

% Foto anggota; jika file belum ada, tampil kotak placeholder
\newcommand{\foto}[1]{%
  \IfFileExists{figures/#1}{\includegraphics[width=3cm,height=4cm,keepaspectratio]{#1}}%
  {\fbox{\parbox[c][4cm][c]{3cm}{\centering\small Foto\\(#1)}}}}

% Sel anggota di cover: nama, NIM, file foto
\newcommand{\anggotacell}[3]{%
  \begin{minipage}[t]{0.45\textwidth}\centering
    \foto{#3}\\[2pt] {\bfseries #1}\\ #2
  \end{minipage}}

% Blok tanda tangan: nama, NIM
\newcommand{\ttd}[2]{%
  \begin{minipage}[t]{0.45\textwidth}\centering
    \vspace{2.5cm}\rule{0.8\linewidth}{0.4pt}\\ #1 (#2)
  \end{minipage}}


\newcommand{\algname}{NAMA\_ALGORITMA}
\newcommand{\repourl}{https://github.com/USERNAME/REPO}
\newcommand{\videourl}{https://youtu.be/XXXXXXXX}
\newcommand{\tahun}{2026}

\newcommand{\anggotaA}{Dita Maheswari}
\newcommand{\nimA}{13523125}
\newcommand{\fotoA}{foto1.jpg}
\newcommand{\anggotaB}{Naufarrel Zhafif Abhista}
\newcommand{\nimB}{13523149}
\newcommand{\fotoB}{foto2.jpg}
\newcommand{\anggotaC}{I Made Wiweka Putera}
\newcommand{\nimC}{13523156}
\newcommand{\fotoC}{foto3.jpg}
\newcommand{\anggotaD}{Hasri Fayadh Muqaffa}
\newcommand{\nimD}{13523160}
\newcommand{\fotoD}{foto4.jpg}

\begin{document}

% ====================================================================
% COVER
% ====================================================================
\begin{titlepage}
  \centering
  {\large\bfseries LAPORAN TUGAS BESAR 2\par}
  {\large\bfseries IF4020 KRIPTOGRAFI\par}
  \vspace{1cm}
  {\Huge\bfseries \algname\par}
  \vspace{0.4cm}
  {\large Perancangan dan Implementasi Block Cipher Modern\par}
  \vspace{1.2cm}

  \vfill
  {\large\bfseries Program Studi Teknik Informatika\\
  Sekolah Teknik Elektro dan Informatika\\
  Institut Teknologi Bandung\\
  \tahun\par}
\end{titlepage}

% ====================================================================
% PERNYATAAN
% ====================================================================
\pagenumbering{roman}
\section*{Pernyataan}
\addcontentsline{toc}{section}{Pernyataan}

\noindent Tugas dan laporan ini disusun sepenuhnya tanpa bantuan kecerdasan buatan
(Generative AI), melainkan hasil pemikiran dan analisis mandiri.

\vspace{1cm}
\noindent
\ttd{\anggotaA}{\nimA}\hfill
\ttd{\anggotaB}{\nimB}

\vspace{1.5cm}
\noindent
\ttd{\anggotaC}{\nimC}\hfill
\ttd{\anggotaD}{\nimD}

\newpage
\tableofcontents
\newpage
\pagenumbering{arabic}

% ====================================================================
% ISI
% ====================================================================
\section{Pendahuluan}

\subsection{Latar Belakang}
% TODO: motivasi perancangan block cipher baru.

\subsection{Review Block Cipher Sejenis}
% TODO: ringkas DES, AES, dll. sebagai pembanding.
\begin{table}[H]
  \centering
  \caption{Perbandingan \algname{} dengan block cipher sejenis}
  \begin{tabular}{llll}
    \toprule
    Parameter & DES & AES-128 & \algname{} \\
    \midrule
    Ukuran blok & 64 bit & 128 bit & \dots{} \\
    Ukuran kunci & 56 bit & 128 bit & \dots{} \\
    Jumlah putaran & 16 & 10 & \dots{} \\
    Struktur & Feistel & SPN & \dots{} \\
    \bottomrule
  \end{tabular}
\end{table}

\subsection{Gagasan dan Pendekatan Utama}
% TODO: ide utama rancangan.

\section{Dasar Teori}
\subsection{Konsep Block Cipher}
\subsection{Jaringan Feistel / SPN}
\subsection{Kotak-S dan Permutasi}
\subsection{Mode Operasi (ECB, CBC, CFB, OFB, CTR)}
\subsection{Padding dan Verifikasi Integritas}
\subsection{Related Works}

\section{Perancangan Block Cipher}
\subsection{Parameter Umum}
\subsection{Alur Enkripsi}
% \begin{figure}[H]
%   \centering
%   \includegraphics[width=0.8\textwidth]{alur-enkripsi.pdf}
%   \caption{Alur enkripsi \algname}\label{fig:enc}
% \end{figure}
\subsection{Alur Dekripsi}
\subsection{Pembentukan \textit{Key Schedule}}
\subsection{\textit{S-Box}}
\subsection{Permutasi / Rotasi}
\subsection{Fungsi Putaran}

\section{Implementasi dan Simulasi}
\subsection{Deskripsi Program}
\subsection{Hasil Uji Enkripsi dan Dekripsi}
\subsubsection{Berkas Teks}
\subsubsection{Berkas Biner}
\begin{table}[H]
  \centering
  \caption{Hasil uji per mode operasi}
  \begin{tabular}{lccc}
    \toprule
    Mode & Berkas teks & Berkas biner & Dekripsi sesuai \\
    \midrule
    ECB & & & \\
    CBC & & & \\
    CFB & & & \\
    OFB & & & \\
    CTR & & & \\
    \bottomrule
  \end{tabular}
\end{table}
\subsection{Penanganan Padding}
\subsection{Verifikasi Integritas Data}

\section{Analisis Keamanan}
\subsection{Avalanche Effect}
\subsection{Entropi}
\subsection{Histogram}
\subsection{Pengujian Tambahan}

\section{Kesimpulan dan Saran Pengembangan}
\subsection{Kesimpulan}
\subsection{Saran Pengembangan (Future Works)}

\begin{thebibliography}{9}
  \bibitem{des} National Institute of Standards and Technology, \textit{FIPS PUB 46-3: Data Encryption Standard (DES)}, 1999.
  \bibitem{aes} National Institute of Standards and Technology, \textit{FIPS PUB 197: Advanced Encryption Standard (AES)}, 2001.
  \bibitem{stinson} D. R. Stinson, \textit{Cryptography: Theory and Practice}, CRC Press.
\end{thebibliography}

\appendix
\section{Lampiran}
\subsection{Repositori Source Code}
\url{\repourl}

\subsection{Video Demo Program}
\url{\videourl}

\subsection{Pembagian Tugas}
\begin{longtable}{@{}p{4cm}p{2.5cm}p{7.5cm}@{}}
  \toprule
  \textbf{Nama} & \textbf{NIM} & \textbf{Kontribusi} \\
  \midrule
  \endhead
  \anggotaA & \nimA & \dots{} \\
  \anggotaB & \nimB & \dots{} \\
  \anggotaC & \nimC & \dots{} \\
  \anggotaD & \nimD & \dots{} \\
  \bottomrule
\end{longtable}


\end{document}
